\section{Metodo Monte Carlo}

Lo studio del modello di Ising mediante simulazioni numeriche si basa su tecniche di campionamento stocastico. 
L’obiettivo è generare configurazioni di spin distribuite secondo la distribuzione di Boltzmann. Poiché il numero totale di configurazioni cresce come $2^N$, un campionamento esaustivo è impraticabile per sistemi macroscopici. 
Si utilizza quindi il metodo Monte Carlo con algoritmo di Metropolis per costruire una catena di Markov che converga alla distribuzione canonica.

\subsection{Algoritmo di Metropolis}

L’algoritmo di Metropolis permette di generare una sequenza di configurazioni $\{C_0, C_1, C_2, \ldots\}$ tale che,
dopo un numero sufficiente di passi,
la distribuzione delle configurazioni sia quella desiderata, ovvero la distribuzione di Boltzmann.
L’algoritmo procede come segue:

\begin{enumerate}
    \item Si sceglie casualmente un sito $i$ del reticolo.
    \item Si propone il flip dello spin $s_i \rightarrow -s_i$.
    \item Si calcola la variazione di energia $\Delta E$ associata al flip.
    \item Il flip viene accettato con probabilità
    \begin{equation}
    P_{\text{acc}} = \min\left(1, e^{-\beta \Delta E}\right).
    \end{equation}
\end{enumerate}

Questo criterio garantisce il rispetto della condizione di bilancio dettagliato:

\begin{equation}
\frac{P(C \rightarrow C')}{P(C' \rightarrow C)} = e^{-\beta (E_{C'} - E_C)},
\end{equation}

assicurando che la distribuzione stazionaria della catena di Markov sia quella canonica.

Un’unità temporale Monte Carlo (MC sweep) corrisponde a $N=L^2$ tentativi di aggiornamento, così che in media ogni spin venga considerato una volta.

\subsection{Dettagli sull'implementazione}

Le simulazioni sono state eseguite su un reticolo quadrato $L \times L$ con condizioni al contorno periodiche, in modo da minimizzare effetti di bordo.

Per ogni temperatura:

\begin{itemize}
    \item È stata effettuata una fase iniziale di termalizzazione, durante la quale le configurazioni non sono state utilizzate per il calcolo delle medie.
    \item Successivamente sono state raccolte configurazioni separate da un numero di 50 sweep per ridurre le correlazioni temporali.
\end{itemize}

Le grandezze fisiche calcolate sono:

\paragraph{Energia media}
\begin{equation}
\langle E \rangle = \frac{1}{N} \sum_{k=1}^{N} E_k
\end{equation}
Dove $E_k$ è l’energia della configurazione $k$ e $N$ è il numero totale di configurazioni campionate.

\paragraph{Magnetizzazione media}
\begin{equation}
\langle |M| \rangle = \frac{1}{N} \sum_{k=1}^{N} |M_k|
\end{equation}
Dove $M_k$ è la magnetizzazione totale della configurazione $k$, si considera il valore assoluto per evitare che le fluttuazioni tra stati con magnetizzazione positiva e negativa si annullino reciprocamente.

Si sono inoltre calcolate il calore specifico $C_V$ e la suscettibilità $\chi$, che sono legate alle fluttuazioni di energia e magnetizzazione rispettivamente Appendice\ref{appendice:calcolo_cv_chi}:
\paragraph{Calore specifico}
\begin{equation}
C_V = \frac{\beta^2}{N} \left( \langle E^2 \rangle - \langle E \rangle^2 \right)
\end{equation}

\paragraph{Suscettibilità}
\begin{equation}
\chi = \frac{\beta}{N} \left( \langle M^2 \rangle - \langle M \rangle^2 \right)
\end{equation}

Inoltre si è stimata la temperatura critica $T_c$ utilizzando il metodo di finite-size scaling, analizzando la posizione del picco della suscettibilità in funzione della dimensione del sistema.

Gli errori statistici sono stati stimati considerando campioni indipendenti e utilizzando l'errore sulla media: ovvero la deviazione standard divisa per la radice del numero di campioni indipendenti.

\section{Risultati della simulazione}

La Figura \ref{fig:simultation_one} mostra l’evoluzione temporale di energia e magnetizzazione per due diverse condizioni iniziali: configurazione ordinata (tutti spin up) e configurazione casuale.

Si osserva che, dopo un certo numero di sweep Monte Carlo, entrambe le evoluzioni convergono agli stessi valori medi, indicando la perdita di memoria della condizione iniziale e il raggiungimento dell’equilibrio termico. Questo comportamento è una verifica numerica dell’ergodicità della dinamica.

\begin{figure}[H]
    \centering
    \begin{subfigure}{0.48\textwidth}
        \centering
        \includegraphics[width=\textwidth]{Mean Energy_tMC.png}
        \caption{Energia media}
        \label{fig:mean_energy}
    \end{subfigure}
    \hfill
    \begin{subfigure}{0.48\textwidth}
        \centering
        \includegraphics[width=\textwidth]{Mean Magnetization_tMC.png}
        \caption{Magnetizzazione media}
        \label{fig:mean_magnetization}
    \end{subfigure}
    \caption{Termalizzazione del sistema per due diverse configurazioni iniziali.}
    \label{fig:simultation_one}
\end{figure}

La Figura \ref{fig:simulation_two} riporta energia media e magnetizzazione media in funzione della temperatura. 


\begin{figure}[H]
    \centering
    \begin{subfigure}{0.48\textwidth}
        \centering
        \includegraphics[width=\textwidth]{fss_mean_energies.png}
        \caption{Energia media}
        \label{fig:mean_energy_vs_T}
    \end{subfigure}
    \hfill
    \begin{subfigure}{0.48\textwidth}
        \centering
        \includegraphics[width=\textwidth]{fss_mean_magnetizations.png}
        \caption{Magnetizzazione media}
        \label{fig:mean_magnetization_vs_T}
    \end{subfigure}
    \caption{Energia media e magnetizzazione media in funzione della temperatura.}
    \label{fig:simulation_two}
\end{figure}

L’energia varia in modo regolare e continuo, mentre la magnetizzazione mostra un comportamento tipico di una transizione di fase: per basse temperature il sistema presenta ordine ferromagnetico con $\langle |M| \rangle \neq 0$, mentre al di sopra di una temperatura critica $T_c$ la magnetizzazione si annulla nel limite termodinamico.

La Figura \ref{fig:simultation_three} mostra il calore specifico e la suscettibilità. Entrambe le quantità presentano un picco in corrispondenza della temperatura critica. 

Il picco del calore specifico è associato alle fluttuazioni energetiche, mentre quello della suscettibilità è legato alle fluttuazioni della magnetizzazione. L’aumento delle fluttuazioni in prossimità di $T_c$ è una firma della criticità del sistema.

Nel limite di sistema infinito, tali grandezze presentano un comportamento non analitico; nel sistema finito simulato si osserva invece un massimo arrotondato, la cui altezza cresce all’aumentare di $L$.

\begin{figure}[H]
    \centering
    \begin{subfigure}{0.48\textwidth}
        \centering
        \includegraphics[width=\textwidth]{fss_heat_capacity.png}
        \caption{Calore specifico}
        \label{fig:heat_capacity}
    \end{subfigure}
    \hfill
    \begin{subfigure}{0.48\textwidth}
        \centering
        \includegraphics[width=\textwidth]{fss_susceptibility.png}
        \caption{Suscettibilità}
        \label{fig:susceptibility}
    \end{subfigure}
    \caption{Capacità termica e suscettibilità in funzione della temperatura.}
    \label{fig:simultation_three}
\end{figure}

\subsection{Stima della temperatura critica e finite-size scaling}

Per stimare la temperatura critica $T_c$ del sistema, si è utilizzato il metodo di finite-size scaling. Esso si basa sull’osservazione che, 
in un sistema finito, le grandezze fisiche non presentano discontinuità ma mostrano un comportamento arrotondato, 
con picchi che si spostano e si accentuano al crescere della dimensione del sistema.

L'equazione di finite-size scaling per la suscettibilità $\chi$ è data da:
\begin{equation}
    T_c^{*}(L) = T_c + a L^{-1/\nu}
\end{equation}
dove $T_c^{*}(L)$ è la temperatura alla quale si osserva il picco della suscettibilità per un sistema di dimensione $L$, $T_c$ è la temperatura critica nel limite termodinamico, $a$ è una costante e $\nu$ è l'esponente critico associato alla lunghezza di correlazione.
Utilizzando i dati ottenuti dalle simulazioni per diverse dimensioni del sistema, è stato possibile effettuare un fit lineare di $T_c^{*}(L)$ in funzione di $L^{-1/\nu}$, con $\nu = 1$ per il modello di Ising bidimensionale.

\begin{figure}[H]
    \centering
    \includegraphics[width=0.6\textwidth]{fss_Tc.png}
    \caption{Stima della temperatura critica $T_c$ mediante finite-size scaling.}
    \label{fig:Tc_fss}
\end{figure}